\section{Colisões com \texttt{pygame.Rect} e o Mini-Jogo}

\begin{frame}{O Poder da Classe \texttt{pygame.Rect}}
    O \textbf{\texttt{pygame.Rect}} é a estrutura mais versátil do Pygame para representar a posição e as dimensões de objetos 2D.

    \vspace{0.3cm}
    \begin{columns}
        \begin{column}{0.5\textwidth}
            \begin{block}{Criação}
                \texttt{r = pygame.Rect(x, y, larg, alt)}
            \end{block}
            \textbf{Propriedades automáticas:}
            \begin{itemize}
                \item \texttt{r.x}, \texttt{r.y}
                \item \texttt{r.width}, \texttt{r.height}
                \item \texttt{r.top}, \texttt{r.bottom}
                \item \texttt{r.left}, \texttt{r.right}
                \item \texttt{r.center}, \texttt{r.centerx}, \texttt{r.centery}
            \end{itemize}
        \end{column}
        \begin{column}{0.5\textwidth}
            \begin{exampleblock}{Bloqueio de Bordas Fácil}
                \begin{itemize}
                    \item \texttt{if r.left < 0: r.left = 0}
                    \item \texttt{if r.right > 800: r.right = 800}
                    \item \texttt{if r.top < 0: r.top = 0}
                    \item \texttt{if r.bottom > 600: r.bottom = 600}
                \end{itemize}
                \small{Não é necessário recalcular largura ou altura manualmente!}
            \end{exampleblock}
        \end{column}
    \end{columns}
\end{frame}

\begin{frame}[fragile]{Detecção de Colisão: \texttt{colliderect()}}
    Para saber se dois objetos se tocaram no espaço 2D, usamos a colisão por caixas alinhadas aos eixos (\textit{AABB - Axis-Aligned Bounding Box}):

    \vspace{0.2cm}
    \begin{lstlisting}[language=Python]
jogador = pygame.Rect(100, 100, 40, 40)
moeda   = pygame.Rect(300, 200, 20, 20)

# O metodo colliderect devolve True se houver sobreposicao!
if jogador.colliderect(moeda):
    pontos += 1
    # Reposiciona a moeda em outro local aleatorio
    moeda.x = random.randint(50, 750)
    moeda.y = random.randint(50, 550)
    \end{lstlisting}

    \vspace{0.2cm}
    \begin{block}{Princípio Matemático AABB}
        Dois retângulos colidem se, e somente se, houver sobreposição simultânea em ambos os eixos $X$ e $Y$:
        \[
            (A_{\text{left}} < B_{\text{right}}) \land (A_{\text{right}} > B_{\text{left}}) \land (A_{\text{top}} < B_{\text{bottom}}) \land (A_{\text{bottom}} > B_{\text{top}})
        \]
    \end{block}
\end{frame}

\begin{frame}[fragile]{Escrevendo Textos na Tela (HUD de Pontuação)}
    Para desenhar pontuações, vidas e mensagens na tela, seguimos três passos:

    \vspace{0.2cm}
    \begin{enumerate}
        \item \textbf{Carregar a Fonte:}
        \begin{lstlisting}[language=Python]
fonte = pygame.font.SysFont("Arial", 24, bold=True)
        \end{lstlisting}
        
        \item \textbf{Renderizar o Texto em uma Surface:}
        \begin{lstlisting}[language=Python]
# Argumentos: (texto, antialiasing_suave, cor_RGB)
img_texto = fonte.render(f"Pontos: {pontos}", True, (255, 255, 255))
        \end{lstlisting}

        \item \textbf{Transferir a Imagem para a Tela (\texttt{blit}):}
        \begin{lstlisting}[language=Python]
# Posiciona a imagem de texto nas coordenadas (x=20, y=20)
tela.blit(img_texto, (20, 20))
        \end{lstlisting}
    \end{enumerate}

    \begin{alertblock}{Atenção}
        A palavra \textbf{blit} vem de \textit{Block Image Transfer} (copiar blocos de pixels de uma superfície para outra).
    \end{alertblock}
\end{frame}
